% titel: Blatt 6, A9: Polare Bindungen und Dipolmoment
% skript: Mortimer Kap. 8/9 (chemie: Dipolmoment, Molekülgeometrie)

Eine Bindung ist polar, wenn die Bindungspartner verschiedene
Elektronegativität haben. Ob das ganze Molekül ein permanentes Dipolmoment
besitzt, ist die Vektorsumme aller Bindungsdipole --- ein symmetrisches
Molekül kann die Einzeldipole kompensieren.

\begin{werkzeugkasten}[polare Bindung vs. Dipolmoment]
\begin{itemize}
  \item \textbf{Polare Bindung:} EN-Differenz $\neq 0$ zwischen den Partnern;
        Elektronen werden zur elektronegativeren Seite gezogen.
  \item \textbf{Permanentes Dipolmoment:} Vektorsumme der Bindungsdipole
        $\neq 0$. Symmetrische Geometrien (linear mit gleichen Enden,
        trigonal-planar mit gleichen Liganden, symmetrisches Tetraeder,
        oktaedrisch) heben die Dipole auf.
  \item Homoatomare Moleküle ($\mathrm{I_2}$, $\mathrm{O_2}$) haben nie ein
        Dipolmoment.
\end{itemize}
\textit{Skript:} Mortimer Kap. 8/9 (Dipol, Molekülgeometrie).
\end{werkzeugkasten}

\begin{achtung}[Darauf achten]
\begin{itemize}
  \item Polare Bindungen \textit{allein} reichen nicht --- erst die
        räumliche Anordnung entscheidet, ob ein Nettodipol bleibt.
  \item $\mathrm{CO_2}$, $\mathrm{CCl_4}$, $\mathrm{CS_2}$ haben polare Bindungen,
        sind aber symmetrisch $\to$ kein Netto-Dipol. $\mathrm{CH_2Cl_2}$ und
        $\mathrm{CHCl_3}$ sind dagegen asymmetrisch $\to$ Dipol.
  \item $\mathrm{SF_4}$ (see-saw) ist asymmetrisch $\to$ Dipol.
\end{itemize}
\end{achtung}
